% =============================================================
%  CHAPTER: Theoretical Framework
%  Thesis: Hallucination Detection in Large Language Models
%  Status: Revised draft, audit corrections applied
% =============================================================

This chapter connects the internal computation of decoder-only Transformer language models with an energy-based view of associative retrieval. The central motivation is that the hidden state entering the feed-forward sublayer can be treated as a query over structured parametric memory. This interpretation is not adopted as a strict architectural identity, but rather as a theoretically motivated modelling perspective supported by prior work on feed-forward layers and associative memory \citep{geva2021transformer, ramsauer2020hopfield}. On that basis, the chapter introduces a Modern Hopfield formulation used to define retrieval distributions and divergence-based indicators for hallucination detection.

% =============================================================

\section{The Transformer Architecture}
\label{sec:transformer}
% ================================================================

The Transformer, introduced by \citet{vaswani2017attention}, replaced the sequential
processing of recurrent networks with a fully parallel, attention-based computation.
It has since become the architectural substrate of virtually all modern language
models. Understanding its internal mechanics is essential to this work: the
hallucination detection framework developed in Chapter~\ref{chap:methods} operates
directly on the per-token hidden states produced at each layer of the forward pass.
Figure~\ref{fig:transformer_block} provides an overview of the complete block
structure; the subsections below describe each component in turn.

% --- Figure 1: full Pre-LN decoder block ----------------------------
\begin{figure}[H]
    \centering
    \includegraphics[width=0.95\linewidth]{figures/fig1_transformer_block.pdf}
    \caption{%
        A single decoder-only Transformer block using Pre-Layer Normalisation
        (RMSNorm) and a SwiGLU feed-forward sublayer, as employed by Qwen2.5 and
        Llama~3. Dashed lines denote residual (skip) connections. The feed-forward
        input $\mathbf{r}_t^{(l)}$ at the output of the second RMSNorm, labelled
        the \emph{probe point}, is the representational signal analysed in
        this work.%
    }
    \label{fig:transformer_block}
\end{figure}
% --------------------------------------------------------------------

% ----------------------------------------------------------------
\subsection{From Recurrent Networks to the Attention Mechanism}
\label{subsec:rnn_to_attention}
% ----------------------------------------------------------------

Before the Transformer, sequence modelling relied primarily on Recurrent Neural Networks (RNNs). These architectures process inputs sequentially: at each time step, the network combines the current token with a hidden state carried over from the previous step, generating an updated state that summarises the sequence history \citep{tunstall2022natural}. However, because each step strictly depends on the previous one, computation cannot be parallelised across positions, making training on long sequences prohibitively slow.

This sequential bottleneck became particularly evident in neural machine translation, where an encoder RNN compressed an entire source sentence into a single, fixed-length vector for the decoder RNN to process. The fundamental weakness of this design was the decoder's limited visibility: it had access only to the \emph{last} encoder hidden state. For long inputs, forcing an arbitrarily long sequence into a fixed-dimensional vector inevitably caused significant information loss \citep{tunstall2022natural}.

To resolve this limitation, the attention mechanism was introduced, granting the decoder direct access to \emph{all} encoder hidden states at every generation step and weighting each by its relevance to the current decoding context \citep{tunstall2022natural}. The Transformer generalised this concept by making attention the primary computational vehicle of the entire architecture, discarding recurrence altogether to enable fully parallelised training.
% ----------------------------------------------------------------
\subsection{Tokens, Embeddings, and Positional Encoding}
\label{subsec:embeddings}
% ----------------------------------------------------------------

A Transformer operates on sequences of discrete units called \emph{tokens}, obtained
by segmenting raw text through a tokeniser. Let $x_1, x_2, \ldots, x_T \in
\mathcal{V}$ denote a sequence of $T$ tokens drawn from vocabulary $\mathcal{V}$.

\paragraph{Token embeddings.}
Each token $x_t$ is mapped to a dense vector via a learned embedding matrix
$\mathbf{E} \in \mathbb{R}^{|\mathcal{V}| \times D}$, where $D$ is the model dimension. This step translates discrete symbols into continuous representations, making the input amenable to gradient-based optimisation. These initial embeddings are context-free, meaning a token receives the same vector regardless of its position or surroundings; dynamic, context-sensitive representations are only introduced later via the self-attention mechanism (Section~\ref{subsec:self_attention}).

\paragraph{Positional encoding.}
Because the Transformer processes all positions in parallel, explicit positional
information must be injected. The original formulation \citep{vaswani2017attention}
used deterministic sinusoidal functions:
\begin{equation}
    \mathbf{p}_t^{(2k)} = \sin\!\left(\frac{t}{10000^{2k/D}}\right), \qquad
    \mathbf{p}_t^{(2k+1)} = \cos\!\left(\frac{t}{10000^{2k/D}}\right),
    \label{eq:sinusoidal_pe}
\end{equation}
where $k \in \{0, \ldots, D/2-1\}$ indexes the embedding dimension. LLMs such as Qwen2.5 \citep{qwen2025qwen25} and Llama~3 \citep{grattafiori2024llama3}, which are the primary model families examined in this work, instead apply \emph{Rotary Positional Encoding} (RoPE) \citep{su2024roformer}. RoPE encodes positional information by rotating query and key vectors prior to their inner product, enabling attention to depend on relative token positions and supporting extrapolation to longer sequences.

Because the models studied here use RoPE, positional information is not added directly to the token embeddings. The initial hidden-state matrix is therefore
\begin{equation}
    \mathbf{H}^{(0)} = \mathbf{E}[x_1, \ldots, x_T] \in \mathbb{R}^{T \times D},
    \label{eq:initial_hidden}
\end{equation}
which serves as the input to the first Transformer block. Unlike the additive sinusoidal scheme of Equation~\eqref{eq:sinusoidal_pe}, RoPE requires no separate positional embedding matrix: positional information is incorporated within each attention layer by rotating the query and key vectors according to their positions before the attention scores are computed.

% ----------------------------------------------------------------
% ----------------------------------------------------------------
\subsection{Self-Attention}
\label{subsec:self_attention}
% ----------------------------------------------------------------

Self-attention is a mechanism that produces a contextualised representation of each token by computing a weighted combination of all token representations in the sequence, where the weights reflect pairwise relevance. To compute these weights, each token simultaneously generates three vectors: a \textit{query} vector representing the information it seeks from the sequence, a \textit{key} vector representing the information it makes available to other tokens, and a \textit{value} vector carrying the information contributed when that token is attended to. These three vectors are obtained via independent learned linear projections of the hidden state matrix $\mathbf{H} \in \mathbb{R}^{T \times D}$:
\begin{equation}
    \mathbf{Q} = \mathbf{H}\mathbf{W}^Q, \qquad
    \mathbf{K} = \mathbf{H}\mathbf{W}^K, \qquad
    \mathbf{V} = \mathbf{H}\mathbf{W}^V,
    \label{eq:qkv_projections}
\end{equation}
where $\mathbf{W}^Q, \mathbf{W}^K \in \mathbb{R}^{D \times d_k}$ and
$\mathbf{W}^V \in \mathbb{R}^{D \times d_v}$ are learned weight matrices. The
relevance of each key to a given query is measured by their scaled inner product;
the resulting scores are normalised via softmax and used to compute a weighted sum
of the value vectors:
\begin{equation}
    \text{Self-Attention}(\mathbf{H}) =
    \text{softmax}\!\left(\frac{\mathbf{Q}\mathbf{K}^\top}{\sqrt{d_k}}\right)\mathbf{V}.
    \label{eq:self_attention}
\end{equation}
The $\sqrt{d_k}$ scaling factor prevents the query-key inner products from growing in magnitude as the dimensionality $d_k$ increases. Without this scaling, the softmax could become excessively peaked, leading to small gradients and less stable optimisation. The output is a context-aware hidden state matrix in which each token representation is a convex combination of the value vectors, weighted according to query-key compatibilities.

Unlike recurrent networks, which accumulate context sequentially through a chain of hidden states, self-attention allows each token to attend directly to any other token in a single parallelisable operation \citep{vaswani2017attention}.

\paragraph{Causal masking.}
In autoregressive LLMs, the hidden representation at position $t$ must depend only on
tokens $x_1, \ldots, x_t$, ensuring that the prediction of token $x_{t+1}$ cannot
access future information. This constraint is enforced during the forward pass by
setting the pre-softmax attention scores to $-\infty$ for all positions $j > t$,
driving the corresponding attention weights to zero. The resulting \emph{causal} or
\emph{masked} self-attention ensures that no token attends to future positions. All
models examined in this work apply causal masking throughout every layer, as
illustrated in Figure~\ref{fig:transformer_block}.

% ----------------------------------------------------------------
\subsection{Multi-Head Attention}
\label{subsec:mha}
% ----------------------------------------------------------------
A single attention head computes one weighted aggregation of value vectors, restricting 
the layer to a single learned relevance pattern. Natural language requires capturing 
diverse, overlapping dependencies within the same layer. \citet{vaswani2017attention} 
addressed this by running $h$ attention heads in parallel, each with its own learned 
projections mapping queries, keys, and values to lower-dimensional subspaces of dimension $d_k = d_v = D/h$. The outputs are concatenated and projected back 
to the model dimension (Figure~\ref{fig:transformer_block}):
\begin{equation}
    \text{MHA}(\mathbf{H}) =
    \text{Concat}(\text{head}_1, \ldots, \text{head}_h)\,\mathbf{W}^O,
    \qquad
    \text{head}_i = \text{Attention}\!\left(
        \mathbf{H}\mathbf{W}^Q_i,\;
        \mathbf{H}\mathbf{W}^K_i,\;
        \mathbf{H}\mathbf{W}^V_i
    \right),
    \label{eq:mha}
\end{equation}
where $\mathbf{W}^O \in \mathbb{R}^{h d_v \times D}$ is the output projection.
Because each head has independent learned projections, different heads can learn to
focus on different relational patterns. Their combination yields a richer
contextualised representation than any single head could produce alone.
% ----------------------------------------------------------------
\subsection{The Feed-Forward Sublayer}
\label{subsec:ffn}
% ----------------------------------------------------------------

Following multi-head attention, each block applies a \emph{position-wise
feed-forward network} (FFN) to every token representation independently. Unlike
attention, which mixes information across positions, the FFN transforms each token
representation $\mathbf{r}_t^{(l)}$ in isolation, applying the same learned function
to every position within a given layer. Its role is therefore complementary:
whereas attention determines \emph{which} information is collected, the FFN
determines \emph{how} that information is transformed and re-encoded at each
position. Because the FFN operates independently on each contextualised token
representation, $\mathbf{r}_t^{(l)}$ (defined precisely in
Section~\ref{subsec:residual_layernorm}) serves as a natural unit of analysis, a
property that this work exploits directly.

In the original Transformer \citep{vaswani2017attention}, the FFN is a two-layer
network with ReLU activation:
\begin{equation}
    \mathrm{FFN}(\mathbf{r})
    =
    \max\!\left(\mathbf{0},\; \mathbf{r}\mathbf{W}_1 + \mathbf{b}_1\right)
    \mathbf{W}_2 + \mathbf{b}_2,
    \label{eq:ffn_original}
\end{equation}
where $\mathbf{W}_1 \in \mathbb{R}^{D \times D_{\mathrm{ff}}}$ and
$\mathbf{W}_2 \in \mathbb{R}^{D_{\mathrm{ff}} \times D}$ are learned weight
matrices, and $D_{\mathrm{ff}}$ is the hidden dimension, typically set to $4D$.

\paragraph{SwiGLU.}
\citet{dauphin2017language} introduced \emph{Gated Linear Units} (GLU) in the
context of gated convolutional language models: an element-wise product between a
sigmoid-gated projection and a linear projection, giving the network an explicit
mechanism for selective information flow. \citet{shazeer2020glu} subsequently
adapted this gating principle to the Transformer FFN and proposed a family of
variants by replacing the sigmoid gate with alternative activation functions. Among these, SwiGLU has become a widely adopted choice in modern LLMs, including Qwen2.5 \citep{qwen2025qwen25}, Llama~3 \citep{grattafiori2024llama3}, Mistral \citep{jiang2023mistral}, and Gemma \citep{gemmateam2024gemma}.
\begin{equation}
    \mathrm{FFN}_{\mathrm{SwiGLU}}(\mathbf{r})
    =
    \left(
        \mathrm{Swish}\!\left(\mathbf{r}\mathbf{W}_1\right)
        \odot
        \mathbf{r}\mathbf{V}
    \right)\mathbf{W}_2,
    \label{eq:swiglu}
\end{equation}
where $\mathbf{W}_1, \mathbf{V} \in \mathbb{R}^{D \times D_{\mathrm{ff}}}$ are the
gate and value projection matrices (following the notation of \citet{shazeer2020glu}),
$\mathbf{W}_2 \in \mathbb{R}^{D_{\mathrm{ff}} \times D}$ is the output projection,
and $\mathrm{Swish}(z) = z \cdot \sigma(z)$ \citep{ramachandran2017searching}.
Figure~\ref{fig:swiglu} details the internal computational graph.
Because SwiGLU uses three weight matrices rather than two, $D_{\mathrm{ff}}$ is
reduced by a factor of $2/3$ relative to the standard $4D$ setting to maintain
comparable parameter counts \citep{shazeer2020glu}.

% --- Figure 2: SwiGLU internal structure ----------------------------
\begin{figure}[H]
    \centering
    \includegraphics[width=0.85\linewidth]{figures/fig2_swiglu_structure.pdf}
    \caption{%
       Internal structure of the SwiGLU feed-forward sublayer
        (Equation~\eqref{eq:swiglu}). The probe point $\mathbf{r}_t^{(l)} \in
        \mathbb{R}^D$ is projected along two parallel branches: a gate branch through
        $\mathbf{W}_1 \in \mathbb{R}^{D \times D_{\mathrm{ff}}}$ followed by the Swish
        activation $\mathrm{Swish}(z) = z \cdot \sigma(z)$, and a value branch through
        $\mathbf{V} \in \mathbb{R}^{D \times D_{\mathrm{ff}}}$, where $D$ is the model
        dimension and $D_{\mathrm{ff}}$ is the intermediate dimension. Their element-wise
        product is projected back to $\mathbb{R}^D$ by $\mathbf{W}_2 \in
        \mathbb{R}^{D_{\mathrm{ff}} \times D}$. The three-matrix parameterisation distinguishes SwiGLU from the original
two-matrix FFN; its relation to the memory view of this work is developed in
Section~\ref{subsec:ffn_memory}.%
    }
    \label{fig:swiglu}
\end{figure}
% --------------------------------------------------------------------

The relation between this parameterisation and associative memory retrieval is
made precise in Section~\ref{subsec:ffn_memory}, where the rows of
$\mathbf{W}_1^{\!\top}$ are identified as the stored patterns of a
layer-specific memory bank and the rows of $\mathbf{W}_2$ as the value
vectors written to the residual stream.
% ----------------------------------------------------------------
\subsection{Residual Connections and Layer Normalisation}
\label{subsec:residual_layernorm}
% ----------------------------------------------------------------

Each Transformer block composes the MHA and FFN sublayers through two structural mechanisms: residual connections and layer normalisation. Together, these mechanisms govern the evolution of the \emph{residual stream}, the sequence of hidden representations that is progressively refined across Transformer layers (Figure~\ref{fig:transformer_block}).

\paragraph{Residual connections.}
A residual (skip) connection adds the sublayer input directly to its output:
\begin{equation}
    \mathbf{y} = \mathbf{x} + f(\mathbf{x}),
    \label{eq:residual}
\end{equation}
where $f(\cdot)$ is the sublayer transformation. This provides a direct gradient
path during backpropagation, mitigating vanishing gradients, and encourages each
sublayer to learn an incremental refinement of its input rather than a complete
transformation from scratch.

\paragraph{Pre-Layer Normalisation.}
Modern LLMs apply normalisation \emph{before} each sublayer (Pre-LN), as shown in
Figure~\ref{fig:transformer_block}. The model comprises $L$ blocks indexed
$l = 0, \ldots, L-1$, matching the layer indexing used throughout
Chapter~\ref{chap:methods}, with the embedding matrix
$\mathbf{H}^{(0)}$ of Equation~\eqref{eq:initial_hidden} as the input to
block $0$. Using $\tilde{\mathbf{H}}^{(l)}$ to denote the
intermediate hidden state after the attention residual of block $l$, the full
block update is:
\begin{align}
    \tilde{\mathbf{H}}^{(l)} &= \mathbf{H}^{(l)}
        + \mathrm{MHA}^{(l)}\!\left(\mathrm{RMSNorm}\!\left(\mathbf{H}^{(l)}\right)\right),
    \label{eq:residual_attn} \\
    \mathbf{H}^{(l+1)} &= \tilde{\mathbf{H}}^{(l)}
        + \mathrm{FFN}^{(l)}\!\left(\mathrm{RMSNorm}\!\left(\tilde{\mathbf{H}}^{(l)}\right)\right).
    \label{eq:residual_ffn}
\end{align}
We fix the naming convention used in the remainder of this work. The symbol
$\mathbf{h}_t^{(l)}$ denotes the \emph{block output} of layer $l$ at position
$t$, i.e.\ the row of $\mathbf{H}^{(l+1)}$ produced by
Equation~\eqref{eq:residual_ffn}; the sequence
$\{\mathbf{h}_t^{(l)}\}_{l=0}^{L-1}$ defines
the residual-stream trajectory of token $t$. Analogously,
$\tilde{\mathbf{h}}_t^{(l)}$ denotes the row of $\tilde{\mathbf{H}}^{(l)}$ at
position $t$, the intermediate state after the attention residual of
Equation~\eqref{eq:residual_attn}. The vector that \emph{enters} the
feed-forward sublayer is given a distinct symbol,
\begin{equation}
    \mathbf{r}_t^{(l)} \;\coloneqq\;
    \mathrm{RMSNorm}\!\left(\tilde{\mathbf{h}}_t^{(l)}\right),
    \label{eq:probe_point}
\end{equation}
the output of the second RMSNorm. The detection framework analyses
$\mathbf{r}_t^{(l)}$ rather than the block output $\mathbf{h}_t^{(l)}$, and we
call it the \emph{probe point} (Figure~\ref{fig:transformer_block}). It is the
argument of the feed-forward function in Equations~\eqref{eq:ffn_original}
and~\eqref{eq:swiglu}, computed immediately before the feed-forward gating
projections.

\paragraph{RMSNorm.}
Both Qwen2.5 and Llama~3 employ Root Mean Square Normalisation (RMSNorm)
\citep{zhang2019rmsnorm}, a simplified variant that omits the mean-centring step
of standard layer normalisation. For a generic row $\mathbf{x} \in \mathbb{R}^D$
of its input matrix,
\begin{equation}
    \mathrm{RMSNorm}\!\left(\mathbf{x}\right)
    =
    \frac{\mathbf{x}}
         {\sqrt{\frac{1}{D}\sum_{d=1}^{D} x_d^2 + \epsilon}}
    \odot \boldsymbol{\gamma},
    \label{eq:rmsnorm}
\end{equation}
where normalisation is applied independently to each row of the input matrix,
$\epsilon > 0$ ensures numerical stability, and $\boldsymbol{\gamma} \in
\mathbb{R}^D$ is a learned scale vector. Removing mean-centring reduces
computational overhead while retaining the stabilising effect on activation
magnitudes.
% ----------------------------------------------------------------
\subsection{Taxonomy of Transformer Architectures}
\label{subsec:transformer_taxonomy}
% ----------------------------------------------------------------

Transformer-based models are commonly grouped into three families according to their
attention pattern \citep{tunstall2022natural}. \emph{Encoder-only} models process
the full input sequence bidirectionally and are primarily used for classification
and language understanding tasks. \emph{Encoder-decoder} models learn a
sequence-to-sequence mapping through a cross-attention bridge between a bidirectional
encoder and a causal decoder, suited to tasks such as machine translation. The
present work is concerned exclusively with \emph{decoder-only} models, which employ
causal masked self-attention to autoregressively generate one token at a time. This family includes Qwen2.5 and Llama 3, which are the architectures examined in the experiments of Chapter~\ref{chap:methods}. The feed-forward inputs $\mathbf{r}_t^{(l)}$ derived from
these states constitute the representational signal that the detection mechanism
analyses.
% =============================================================
% =============================================================
% ================================================================
\section{Internal Representations of LLMs}
\label{sec:internal_representations}
% ================================================================
Using the notation established in Section~\ref{subsec:residual_layernorm}, parametric memory retrieval operates on the feed-forward input $\mathbf{r}_t^{(l)}$ of Equation~\eqref{eq:probe_point}. We now make this interpretation of the feed-forward sublayer precise.
% ----------------------------------------------------------------
% ----------------------------------------------------------------
\subsection{Feed-Forward Layers as Parametric Memory}
\label{subsec:ffn_memory}
% ----------------------------------------------------------------

Each feed-forward sublayer processes token representations independently, applying
the same learned transformation to every position. This position-wise independence
means that the probe point $\mathbf{r}_t^{(l)}$ acts on the sublayer's parameters
without interacting with any other token, making it a natural unit for studying
how individual representations interact with learned knowledge.

\citet{geva2021transformer} formalised this observation by showing that the
feed-forward sublayer can be written as an unnormalised key-value memory lookup.
Letting $\mathbf{r} \in \mathbb{R}^D$ denote the token representation entering the
sublayer (the probe point $\mathbf{r}_t^{(l)}$ of
Equation~\eqref{eq:probe_point}), the computation takes the form:
\begin{equation}
    \mathrm{FFN}^{(l)}(\mathbf{r})
    =
    g\!\left(\mathbf{r}\,{\mathbf{K}_{\mathrm{FFN}}^{(l)}}^{\!\top}\right)
    \mathbf{V}_{\mathrm{FFN}}^{(l)},
    \label{eq:ffn_kv_view}
\end{equation}
where $\mathbf{K}_{\mathrm{FFN}}^{(l)} \in \mathbb{R}^{K_l \times D}$ is the first
weight matrix, whose $K_l$ rows act as key vectors,
$\mathbf{V}_{\mathrm{FFN}}^{(l)} \in \mathbb{R}^{K_l \times D}$ is the second
weight matrix, whose rows act as value vectors, and $g(\cdot)$ is an element-wise
nonlinearity such as ReLU. The subscript distinguishes these matrices from the
query, key, and value projections of the attention sublayer
(Equation~\eqref{eq:qkv_projections}) and from the SwiGLU value projection
$\mathbf{V}$ of Equation~\eqref{eq:swiglu}, which are unrelated objects. The inner
products $\mathbf{r}\,{\mathbf{K}_{\mathrm{FFN}}^{(l)}}^{\!\top}$ measure the
alignment between the input and each key, the nonlinearity $g$ produces a
coefficient for each memory slot, and $\mathbf{V}_{\mathrm{FFN}}^{(l)}$ assembles
the output as a weighted combination of value vectors. The distinction from a
normalised memory is that the canonical Transformer does not apply softmax to
these coefficients, making the retrieval unnormalised \citep{geva2021transformer}.

This interpretation is supported empirically by \citet{geva2021transformer}, who
found that individual key vectors often correspond to recurring input patterns.
Keys in lower layers tend to respond to shallow lexical or syntactic features,
whereas those in higher layers are associated with more semantic patterns. The
corresponding value vectors influence the model's predictions in ways consistent
with these activations, providing evidence that feed-forward layers function as
distributed repositories of learned patterns.
% ----------------------------------------------------------------
\subsection{Gated Feed-Forward Layers and the Memory Bank}
\label{subsec:gated_ffn_memory}
% ----------------------------------------------------------------
The models examined in this work, Qwen2.5 and Llama~3, use the SwiGLU variant
described in Section~\ref{subsec:ffn} rather than the ReLU-based FFN of
Equation~\eqref{eq:ffn_kv_view}. In generic form, a gated feed-forward layer
computes two parallel projections of its input and combines them
multiplicatively:
\begin{equation}
    \mathrm{FFN}^{(l)}(\mathbf{r})
    =
    \Bigl(
        \phi\!\left(\mathbf{r}\,{\mathbf{A}^{(l)}}^{\!\top}\right)
        \odot
        \mathbf{r}\,{\mathbf{B}^{(l)}}^{\!\top}
    \Bigr)\mathbf{C}^{(l)},
    \label{eq:gated_ffn_generic}
\end{equation}
where $\mathbf{A}^{(l)}, \mathbf{B}^{(l)} \in \mathbb{R}^{K_l \times D}$,
$\mathbf{C}^{(l)} \in \mathbb{R}^{K_l \times D}$, $\phi(\cdot)$ is an
element-wise nonlinearity, and $\odot$ denotes element-wise multiplication.
Here $K_l$ denotes the number of intermediate units of layer $l$; for the
SwiGLU parameterisation of Equation~\eqref{eq:swiglu} this number equals the
intermediate dimension, $K_l = D_{\mathrm{ff}}$, an identification assumed
throughout this work. Bias terms are omitted, consistent with the formulation
of \citet{shazeer2020glu} and with the parameterisation of the models studied
here, whose feed-forward projections carry no bias. In the SwiGLU case of
Equation~\eqref{eq:swiglu}, written there in the row-vector convention
$\mathbf{r}\mathbf{W}_1$ with $\mathbf{W}_1, \mathbf{V} \in
\mathbb{R}^{D \times D_{\mathrm{ff}}}$, the row-indexed matrices of the gated
form are obtained by transposing the gate and value projections, whose layer
superscript, suppressed in Equation~\eqref{eq:swiglu}, is restored here,
\begin{equation}
    \mathbf{A}^{(l)} = {\mathbf{W}_1^{(l)}}^{\!\top}, \qquad
    \mathbf{B}^{(l)} = {\mathbf{V}^{(l)}}^{\!\top}, \qquad
    \mathbf{C}^{(l)} = \mathbf{W}_2^{(l)},
    \label{eq:gated_identification}
\end{equation}
so that each \emph{row} of $\mathbf{A}^{(l)}$ is one gate pattern in
$\mathbb{R}^D$ (see Figure~\ref{fig:swiglu}). The transposition involves no
change in the underlying parameters; it only reorganises the matrices so that
stored patterns are indexed by rows, which is the convention required for the
memory bank defined below.

The gating structure has a consequence that the ReLU case does not share. In
the architecture analysed by \citet{geva2021transformer}, the coefficient of
each intermediate unit is a function of a single inner product between the
input and one row of the first weight matrix, which makes that matrix the
unique structural candidate for the role of stored keys. In
Equation~\eqref{eq:gated_ffn_generic}, the coefficient of unit $i$ is instead
the product $\phi(\mathbf{r}\,\mathbf{a}_i^{(l)})\,
(\mathbf{r}\,\mathbf{b}_i^{(l)})$, where
$\mathbf{a}_i^{(l)}, \mathbf{b}_i^{(l)} \in \mathbb{R}^{D}$ are the vectors
whose transposes form the $i$-th rows of $\mathbf{A}^{(l)}$ and
$\mathbf{B}^{(l)}$. The retrieval score of each unit therefore depends on two
projections of the input simultaneously, and no single weight matrix occupies
the position that the key matrix occupies in the ReLU case. Any identification
of a fixed memory bank with a specific weight matrix is consequently a
modelling choice rather than a structural theorem, a point consistent with
recent analyses of associative memory in gated feed-forward blocks
\citep{guptalinearly}. The key-value view of \citet{geva2021transformer}
applies directly to the ReLU architecture and requires a deliberate extension
when applied to gated ones.

In this thesis, the memory bank of layer $l$ is defined as the collection of
row vectors of $\mathbf{A}^{(l)}$:
\begin{equation}
    \mathbf{M}^{(l)} \coloneqq \mathbf{A}^{(l)}
    =
    \begin{bmatrix}
        {\mathbf{m}_1^{(l)}}^{\!\top} \\
        \vdots \\
        {\mathbf{m}_{K_l}^{(l)}}^{\!\top}
    \end{bmatrix}
    \in \mathbb{R}^{K_l \times D}.
    \label{eq:memory_bank_definition}
\end{equation}
This choice is motivated by the position of $\mathbf{A}^{(l)}$ within the
gating computation. The gate branch is the only branch passed through the
nonlinearity, so the term
$\phi(\mathbf{r}\,{\mathbf{A}^{(l)}}^{\!\top})$ determines which intermediate
units are selectively activated, in direct parallel with the role of
$g(\mathbf{r}\,{\mathbf{K}_{\mathrm{FFN}}^{(l)}}^{\!\top})$ in the formulation
of \citet{geva2021transformer}. The branch through $\mathbf{B}^{(l)}$ enters the
computation linearly and modulates the magnitude of each activated unit rather
than its selection, while $\mathbf{C}^{(l)}$ writes the result back to the
residual stream. Both matrices remain essential to the full computation, but
neither serves the pattern-matching role that defines a key in the sense of
\citet{geva2021transformer}.

With this definition, the probe point $\mathbf{r}_t^{(l)}$ entering the
feed-forward sublayer is treated as a query over the memory bank
$\mathbf{M}^{(l)}$: the inner products
$\mathbf{r}_t^{(l)}{\mathbf{M}^{(l)}}^{\!\top}$ measure the alignment of the
token's representation with each stored pattern. Section~\ref{sec:hopfield}
formalises this retrieval operation within a Modern Hopfield Network framework,
which provides an energy function and a closed-form retrieval distribution over
the memory bank.
% ================================================================
% ================================================================
\section{Modern Hopfield Networks}
\label{sec:hopfield}
% ================================================================

Section~\ref{subsec:gated_ffn_memory} established the memory bank
$\mathbf{M}^{(l)}$ and interpreted the feed-forward input
$\mathbf{r}_t^{(l)}$ as a query whose inner products with the stored patterns
measure alignment. That interpretation remains informal until a framework
specifies what retrieval means quantitatively. Modern Hopfield networks supply
exactly this framework: given a query state and a bank of stored patterns, they
define an energy function whose minima correspond to retrieved patterns, an
update rule that descends that energy, and a closed-form softmax distribution
over the bank. The energy and the distribution are the two objects this thesis
extracts at every layer and token position, and their evolution between the
prefill and generation phases is the raw material of the detection signal
developed in Section~\ref{sec:divergences}.

% ----------------------------------------------------------------
\subsection{From Classical to Modern Associative Retrieval}
\label{subsec:classical_to_modern_hopfield}
% ----------------------------------------------------------------

The classical Hopfield network \citep{hopfield1982neural} stores binary
patterns in a symmetric weight matrix
$\mathbf{W} = \sum_{i=1}^{K} \mathbf{m}_i\mathbf{m}_i^\top
\in \mathbb{R}^{D \times D}$ constructed by Hebbian learning, and retrieves
them by iterating a sign update that decreases the quadratic energy
\begin{equation}
    E_{\mathrm{cl}}(\mathbf{q})
    =
    -\tfrac{1}{2}\mathbf{q}^{\top}\mathbf{W}\mathbf{q},
    \qquad \mathbf{q} \in \{-1,+1\}^D.
    \label{eq:classical_hopfield_energy}
\end{equation}
Two restrictions make this formulation unusable here. The states are binary,
whereas Transformer hidden states are continuous vectors, and the storage
capacity grows only linearly in the dimension $D$, whereas the feed-forward
banks of Section~\ref{subsec:gated_ffn_memory} contain $K_l > D$ patterns,
well in excess of that linear capacity.

\citet{ramsauer2020hopfield} resolve both restrictions with a modern Hopfield
network over continuous states whose capacity, under separation conditions on
the stored patterns, grows exponentially in $D$. In this section and for the
remainder of the chapter, single states such as the query are written as
column vectors; the row-form inner products
$\mathbf{r}_t^{(l)}{\mathbf{M}^{(l)}}^{\!\top}$ of
Section~\ref{subsec:gated_ffn_memory} accordingly reappear below in the
transposed form ${\mathbf{m}_i^{(l)}}^{\!\top}\mathbf{r}_t^{(l)}$. Let
\begin{equation}
    \mathbf{M}
    =
    \begin{bmatrix}
        \mathbf{m}_1^\top\\
        \vdots\\
        \mathbf{m}_K^\top
    \end{bmatrix}
    \in \mathbb{R}^{K \times D}
    \label{eq:generic_memory_bank}
\end{equation}
be a memory bank whose $K$ rows are stored patterns, and let
$\mathbf{q} \in \mathbb{R}^D$ be a continuous query. The modern Hopfield
energy is
\begin{equation}
    E(\mathbf{q};\,\mathbf{M},\beta)
    =
    -\frac{1}{\beta}\log\sum_{i=1}^{K}
    \exp\!\bigl(\beta\,\mathbf{m}_i^{\top}\mathbf{q}\bigr)
    +\frac{1}{2}\|\mathbf{q}\|_2^{2}
    + c_0(\mathbf{M},\beta,K),
    \label{eq:modern_hopfield_energy}
\end{equation}
where $\beta > 0$ is an inverse temperature and
$c_0(\mathbf{M},\beta,K) = \beta^{-1}\log K +
\frac{1}{2}\max_i\|\mathbf{m}_i\|_2^2$ is an additive constant that depends on
the bank but not on the query. The first term rewards alignment with the
stored patterns through a log-sum-exp over inner products; the quadratic term
penalises queries of large norm and guarantees that the energy is bounded
below. Since $c_0$ is constant in $\mathbf{q}$, it cancels in every comparison
between energies of different queries against the same bank and is omitted in
the computations of this thesis.

The energy is minimised by the update rule \citep{ramsauer2020hopfield}
\begin{equation}
    \mathbf{q}^{+}
    =
    \mathbf{M}^{\top}\,\mathrm{softmax}\!\bigl(\beta\,\mathbf{M}\mathbf{q}\bigr),
    \label{eq:modern_hopfield_update}
\end{equation}
which replaces the query by a convex combination of the stored patterns,
weighted by a softmax over their inner products with the current state.
\citet{ramsauer2020hopfield} prove global convergence of this iteration and
characterise its fixed points. When no pattern dominates, the iteration
converges to a global fixed point near the arithmetic mean of all patterns,
and the softmax is near uniform. When a correlated subset of patterns jointly
attracts the query, the iteration settles into a metastable state that
averages over that subset. When the query lies well within the basin of a
single well-separated pattern, the iteration recovers that pattern alone.
This trichotomy is the theoretical basis for reading the concentration of the
softmax as a retrieval-quality signal, an interpretation made operational in
Section~\ref{subsec:retrieval_distribution_llm}.

\paragraph{Connection to the attention mechanism.}
\citet{ramsauer2020hopfield} show that the update rule of
Equation~\eqref{eq:modern_hopfield_update} coincides with scaled dot-product
attention under the identification $\beta = 1/\sqrt{d_k}$, with the stored
patterns serving simultaneously as keys and values:
\begin{equation}
    \mathbf{q}^{+}
    =
    \mathrm{Attention}\!\left(\mathbf{q},\,\mathbf{M},\,\mathbf{M}\right)
    =
    \mathbf{M}^{\top}\,\mathrm{softmax}\!\!\left(
        \frac{\mathbf{M}\mathbf{q}}{\sqrt{d_k}}
    \right).
    \label{eq:mhn_attention_equiv}
\end{equation}
This is the attention operation of \citet{vaswani2017attention} with
$\mathbf{K} = \mathbf{V} = \mathbf{M}$. It is not self-attention, where queries, keys, and
values arise from separate learned projections of the same sequence. The
correspondence nonetheless matters for the present argument: the Transformer
and the modern Hopfield network share the same core retrieval operation, so
applying Hopfield machinery to the internals of a Transformer is not an
external imposition but a reading of the architecture in its own native terms.

\paragraph{Sparse extensions.}
Later work generalises the dense softmax retrieval map to sparse alternatives
such as sparsemax and $\alpha$-entmax, which strengthen convergence to
single-pattern fixed points \citep{martins2023sparse}. This thesis retains the
dense softmax formulation deliberately: it produces a strictly positive
probability distribution over all stored patterns, which is the object
required by the entropy and divergence analysis of
Section~\ref{sec:divergences}, and its strict positivity guarantees that the
Kullback-Leibler divergence introduced there is always finite.

% ----------------------------------------------------------------
\subsection{Retrieval Distribution Induced by an LLM Hidden State}
\label{subsec:retrieval_distribution_llm}
% ----------------------------------------------------------------

The detection problem of this thesis does not require running the Hopfield
iteration to convergence. The model itself advances the computation; what the
analysis needs is a snapshot of how a hidden state relates to the memory bank
at the moment it enters the feed-forward sublayer. A single application of the
softmax map provides that snapshot.

Instantiating the generic framework of
Section~\ref{subsec:classical_to_modern_hopfield} with the Transformer
quantities of Sections~\ref{sec:internal_representations} and
\ref{subsec:gated_ffn_memory}, the query is the feed-forward input at layer
$l$ and position $t$,
\begin{equation}
    \mathbf{q} = \mathbf{r}_t^{(l)},
    \label{eq:query_identification}
\end{equation}
and the bank is the layer-specific matrix $\mathbf{M}^{(l)}$ of
Equation~\eqref{eq:memory_bank_definition}. The induced retrieval distribution
is
\begin{equation}
    p_i^{(l,t)}
    =
    \frac{
        \exp\!\bigl(\beta\,{\mathbf{m}_i^{(l)}}^{\!\top}\mathbf{r}_t^{(l)}\bigr)
    }{
        \sum_{j=1}^{K_l}
        \exp\!\bigl(\beta\,{\mathbf{m}_j^{(l)}}^{\!\top}\mathbf{r}_t^{(l)}\bigr)
    },
    \qquad i = 1, \ldots, K_l.
    \label{eq:retrieval_distribution}
\end{equation}
This distribution converts a token-level hidden state into a probability
vector over the $K_l$ stored patterns of layer $l$. The retrieved prototype,
that is, the pattern the hidden state most nearly recovers, is the expectation
under this distribution:
\begin{equation}
    \bar{\mathbf{m}}_t^{(l)}
    =
    \sum_{i=1}^{K_l} p_i^{(l,t)}\,\mathbf{m}_i^{(l)}.
    \label{eq:retrieved_pattern}
\end{equation}
The retrieved prototype is not itself a detection feature; it is stated here
because it makes explicit what the distribution retrieves, and it grounds the
interpretive reading of the fixed-point trichotomy that follows.

Equation~\eqref{eq:retrieval_distribution} is the central bridge between the
Transformer and the Hopfield formalism, and every detection feature in this
thesis is a functional of it. In the operational form of
Chapter~\ref{chap:methods}, the inner products entering this softmax are
evaluated with the stored patterns rescaled to unit norm; the construction
and the resulting distinction between the model's native pre-activations and
the derived scores are stated there. Its shape admits a direct reading through the
fixed-point trichotomy of \citet{ramsauer2020hopfield}. A near-uniform
distribution corresponds to the global averaging regime, in which the hidden
state engages no specific stored knowledge. A distribution concentrated on a
small subset corresponds to a metastable state. A distribution sharply peaked
on one pattern indicates that the query sits inside the basin of attraction of
a single memory. A natural scalar summary of this concentration is the
retrieval entropy
\begin{equation}
    H^{(l,t)}
    =
    -\sum_{i=1}^{K_l} p_i^{(l,t)} \log p_i^{(l,t)}.
    \label{eq:retrieval_entropy}
\end{equation}
High entropy indicates diffuse retrieval; low entropy indicates that the
hidden state is sharply aligned with few stored patterns.

The hypothesis motivating this work can now be stated in these terms. During
question prefill, the model processes grounded input and its hidden states
induce retrieval distributions with some characteristic structure per layer.
During answer generation, a model producing content supported by its
parametric knowledge should induce distributions whose structure remains
consistent with that reference, whereas a model hallucinating should drift
toward retrieval patterns that diverge from it. The direction of this drift
is not assumed. The fixed-point trichotomy admits a second failure mode:
generation may settle deep inside the basin of a coherent but factually
wrong memory, producing low energy and sharp concentration, a confidently
wrong retrieval whose defect lies in content rather than in stability. What
the hypothesis asserts is the comparative
structure of the signal, namely that the relationship between prefill and
generation retrieval differs systematically between grounded and hallucinated
answers; the sign and depth profile of that difference are empirical
questions operationalised in Chapter~\ref{chap:methods}. Testing this hypothesis
requires quantitative measures of dissimilarity between distributions, which
the next section provides.

% ================================================================
\section{Information-Theoretic Divergences}
\label{sec:divergences}
% ================================================================

The retrieval distribution of Equation~\eqref{eq:retrieval_distribution}
describes a single hidden state at one layer and one position. The detection
hypothesis stated at the end of
Section~\ref{subsec:retrieval_distribution_llm} is comparative: it concerns
the relationship between distributions induced during question prefill and
distributions induced during answer generation. This section introduces the
quantities used to make that comparison precise. Two of them, entropy and its
normalised form, summarise a single distribution; the remaining ones, the
Kullback-Leibler, Jensen-Shannon, and Hellinger measures, quantify the
dissimilarity between two distributions. Throughout, $P = (p_1, \ldots, p_K)$
and $Q = (q_1, \ldots, q_K)$ denote probability distributions over $K$
outcomes; these symbols are distribution labels and are distinct from the
query and key matrices of Section~\ref{subsec:self_attention}.

% ----------------------------------------------------------------
\subsection{Shannon Entropy}
\label{subsec:shannon_entropy}
% ----------------------------------------------------------------

The Shannon entropy of a discrete distribution $P$
\citep{shannon1948mathematical} measures the average uncertainty of drawing
from it:
\begin{equation}
    H(P) = -\sum_{i=1}^{K} p_i \log p_i.
    \label{eq:shannon_entropy}
\end{equation}
It attains its maximum $\log K$ at the uniform distribution and its minimum
$0$ at a point mass. The retrieval entropy $H^{(l,t)}$ of
Equation~\eqref{eq:retrieval_entropy} is the instance of
Equation~\eqref{eq:shannon_entropy} with $P = \{p_i^{(l,t)}\}_{i=1}^{K_l}$.

Because the bank size $K_l$ differs across model architectures, and is not
constrained by the framework to be uniform across layers, raw entropies are
not directly comparable between banks of different size. The analysis therefore
uses the \emph{normalised entropy}
\begin{equation}
    \bar{H}(P) = \frac{H(P)}{\log K},
    \label{eq:normalised_entropy}
\end{equation}
which maps the range of $H$ to $[0,1]$ and makes per-layer and cross-model
comparisons dimensionless. A value of $\bar{H} \approx 1$ indicates
near-uniform retrieval, the global averaging regime of the Hopfield
trichotomy; $\bar{H} \approx 0$ indicates retrieval concentrated on a single
pattern.

% ----------------------------------------------------------------
\subsection{Kullback-Leibler Divergence}
\label{subsec:kl_divergence}
% ----------------------------------------------------------------

Comparing two distributions begins with the Kullback-Leibler (KL) divergence
\citep{kullback1951information}, which measures the expected log-likelihood
ratio of $P$ over $Q$ under $P$:
\begin{equation}
    D_{\mathrm{KL}}(P \| Q) = \sum_{i=1}^{K} p_i \log \frac{p_i}{q_i}.
    \label{eq:kl_divergence}
\end{equation}
By Gibbs' inequality, $D_{\mathrm{KL}}(P \| Q) \geq 0$, with equality if and
only if $P = Q$. The divergence is asymmetric in general, and it is finite
only when the support of $P$ is contained in the support of $Q$. In the
present application both distributions are softmax outputs over the full
memory bank, so all entries are strictly positive and the divergence is
always well defined. The KL divergence is not used directly as a feature; it
is retained as the building block from which the symmetric measure of the
next subsection is constructed.

% ----------------------------------------------------------------
\subsection{Jensen-Shannon Divergence}
\label{subsec:js_divergence}
% ----------------------------------------------------------------

A feature intended to compare prefill and generation phases should not depend
on which phase is designated as reference. The Jensen-Shannon divergence
(JSD) \citep{lin1991divergence} achieves this by symmetrising the KL
divergence through the mixture $M = \tfrac{1}{2}(P + Q)$:
\begin{equation}
    D_{\mathrm{JS}}(P \| Q)
    =
    \frac{1}{2}\,D_{\mathrm{KL}}(P \| M)
    +
    \frac{1}{2}\,D_{\mathrm{KL}}(Q \| M).
    \label{eq:js_divergence}
\end{equation}
An equivalent expression in terms of entropies makes the connection to
Equation~\eqref{eq:shannon_entropy} explicit:
\begin{equation}
    D_{\mathrm{JS}}(P \| Q)
    =
    H(M) - \frac{1}{2}H(P) - \frac{1}{2}H(Q).
    \label{eq:js_entropy_form}
\end{equation}
The JSD is symmetric and, with natural logarithms, satisfies
$0 \leq D_{\mathrm{JS}}(P \| Q) \leq \log 2$ regardless of the supports of
$P$ and $Q$. Boundedness matters in practice: a feature fed to a downstream
classifier should not be dominated by a handful of unbounded values. These
properties make the JSD the primary divergence of this work for comparing
retrieval distributions across the prefill and generation phases.

% ----------------------------------------------------------------
\subsection{Hellinger Distance}
\label{subsec:hellinger}
% ----------------------------------------------------------------

A single divergence captures one notion of dissimilarity, and a probe trained
on its output inherits its blind spots. The squared Hellinger distance
\citep{hellinger1909neue} is included as a geometrically distinct companion
to the JSD:
\begin{equation}
    H^2_{\mathrm{Hel}}(P, Q)
    =
    \frac{1}{2}\sum_{i=1}^{K}
    \!\left(\sqrt{p_i} - \sqrt{q_i}\right)^{\!2}.
    \label{eq:hellinger_squared}
\end{equation}
It satisfies $0 \leq H^2_{\mathrm{Hel}}(P, Q) \leq 1$, with
$H^2_{\mathrm{Hel}} = 0$ if and only if $P = Q$ and
$H^2_{\mathrm{Hel}} = 1$ if and only if $P$ and $Q$ have disjoint supports.
It is symmetric and bounded without any support condition, and it operates on
square-root probabilities rather than log-ratios, which keeps its computation
numerically stable near small probabilities.

Its relationship to the Bhattacharyya coefficient
$\mathrm{BC}(P, Q) = \sum_{i=1}^{K} \sqrt{p_i q_i}$ is direct:
\begin{equation}
    H^2_{\mathrm{Hel}}(P, Q) = 1 - \mathrm{BC}(P, Q),
    \label{eq:hellinger_bhattacharyya}
\end{equation}
so the squared Hellinger distance measures how far the two distributions are
from sharing identical square-root profiles. Where the JSD is built from
log-ratios and weighs discrepancies through the mixture, the Hellinger
distance compares probability amplitudes pointwise, giving it complementary
sensitivity to differences in the tails. The two measures therefore encode
related but non-redundant views of the same distributional shift.

In this work, $D_{\mathrm{JS}}$ and $H^2_{\mathrm{Hel}}$ are computed at each
layer $l$ by comparing the retrieval distribution of each generated token
against a prefill reference distribution from the same layer, and the
resulting per-token scalars are averaged across the generated sequence to
yield per-layer divergence scores. Together with the energy of
Equation~\eqref{eq:modern_hopfield_energy} and the normalised entropy of
Equation~\eqref{eq:normalised_entropy}, these scores constitute the layerwise
features of the detection method. The precise definition of the prefill
reference and the construction of the final feature vector are formalised in
Chapter~\ref{chap:methods}.
% =============================================================
%  END OF CHAPTER
% =============================================================